# Introduction

The status quo of generative AI safety treats prompt injection and
jailbreaks as semantic problems, utilizing auxiliary linguistic scanners
to evaluate textual intent. To an adversary, natural language represents
an infinitely flexible obfuscation surface; an execution payload can be
dynamically masked using unique idioms, structural formatting fragments,
or polymorphic syntax that completely evades keyword or semantic
matching.

HGFW introduces a systems-engineering paradigm shift: we treat context
tracking not as a text-parsing exercise, but as a deterministic
geometric trajectory through a continuous, high-dimensional vector
space. An inbound user query does not merely present words; it
introduces a mathematical tensor that pulls the model's internal
computational state toward specific coordinate bounds. HGFW establishes
a default-deny boundary directly at the protocol gate, verifying the
structural capability of a vector payload before it is ever permitted to
reach the primary Transformer layers.

# Technical Architecture & Mathematical Formalization

The operational pipeline of the HGFW engine executes sequentially across
three linear algebra layers during the embedding pre-processing phase,
preceding the core transformer layers.

## Embedding Anisotropy Mitigation (The Cone Correction)

Raw hidden-state vectors extracted from transformer-based text embedders
inherently suffer from severe anisotropy; sentence representation
vectors cluster tightly into a highly constrained, narrow directional
"cone" within the hyperdimensional manifold ($\mathbb{R}^d$). To strip
away this uniform, shared background noise and maximize semantic
variance resolution, HGFW enforces an upfront mean-centering
transformation:

$$\begin{equation}
\mu = \frac{1}{M}\sum_{i=1}^{M} A_i
\end{equation}$$ $$\begin{equation}
x = x_{\text{raw}} - \mu
\end{equation}$$

Where $A \in \mathbb{R}^{M \times d}$ represents an authenticated
reference corpus of $M$ safe, domain-specific user interactions,
$\mu \in \mathbb{R}^d$ is the calculated global mean vector, and $x$
represents the corrected, zero-centered inbound payload vector.

## Zero-Trust Subspace Projection

Rather than attempting to model or anticipate an infinite, unpredictable
array of malicious directions, HGFW implements a default-deny **Valid
Context Subspace ($V$)**. We construct this whitelist manifold by
executing Singular Value Decomposition (SVD) on the centered training
matrix $A_{\text{centered}}$:

$$\begin{equation}
A_{\text{centered}} = U \Sigma V^T
\end{equation}$$

Where $V^T \in \mathbb{R}^{d \times d}$ holds the orthonormal basis
vectors sorted in descending order by informational magnitude
(eigenvalues). We extract the top $k$ primary dimensions to define our
approved subspace, yielding the truncated matrix
$V_k \in \mathbb{R}^{d \times k}$. We then compile a static, symmetric
orthogonal projection matrix ($P_V$) configured strictly for our
authorized domain coordinates:

$$\begin{equation}
P_V = V_k V_k^T
\end{equation}$$

For every incoming centered payload vector $x$, the non-whitelisted
structural attributes are isolated into the **orthogonal residual vector
($e$)**:

$$\begin{equation}
e = x - P_V x
\end{equation}$$ $$\begin{equation}
\|e\| = \sqrt{\sum_{j=1}^{d} e_j^2}
\end{equation}$$

## Signed Structural Variance Tracking

To close the rate-limiting and session-reset loopholes where an
adversarial botnet interleaves malicious probes between large blocks of
benign traffic, HGFW tracks the directional velocity of the residual
trace over time.

Benign human context transitions smoothly across subsequent turns.
Conversely, strategic adversarial state-switching forces a violent
geometric whiplash as the vector sequence moves abruptly out of the
whitelisted subspace. The engine computes the signed growth delta
($\Delta_{\text{growth}}$) turn-over-turn:

$$\begin{equation}
\Delta_{\text{growth}} = \|e_t\| - \|e_{t-1}\|
\end{equation}$$

If $\Delta_{\text{growth}}$ exceeds a calibrated threshold window
($\text{margin}$), a **Structural Variance Trigger** activates, dropping
the execution gate instantly. By tracking the *signed directional
growth* ($\Delta_{\text{growth}} > 0$) rather than the absolute scalar
difference, the firewall prevents false-positive feedback loops when a
legitimate user transitions safely back down toward whitelisted
coordinates (negative growth).

# Experimental Evaluation & Empirical Telemetry

To verify the theoretical framework, HGFW was integrated directly into a
production-grade 768-dimensional language embedder
(`bert-base-uncased`). The whitelist manifold was calibrated on a sample
corpus mapping software engineering queries, with hyperparameters tuned
to $k=3$ (whitelisted dimensions) and $\text{margin}=0.50$.

A multi-turn adversarial exploit and camouflage sequence was fed into
the engine. The resulting raw empirical data telemetry is documented in
Table 1:

   Turn ($t$)    Payload Class    $\|e\|$   $\Delta_{\text{growth}}$    Status
  ------------ ----------------- --------- -------------------------- -----------
       1          Benign Init     1.2988             *Base*            **PASS**
       2         Benign Domain    1.4613            +0.1625            **PASS**
       3        Admin Override    2.8794          **+1.4181**          **BLOCK**
       4        Camouflage Bait   2.7148            -0.1646            **BLOCK**
       5         Exploit Exec     2.0650          **+0.6037**          **BLOCK**

  : HGFW Empirical Telemetry Trace

## Telemetry Post-Mortem & Lexical Overlap Analysis

The empirical data exposes a crucial behavioral signature. On Turn 3,
the prompt injection causes an immediate, massive directional escape
from the SVD whitelist space, resulting in a growth delta of
**+1.4181**, which safely triggers the variance blocker.

On Turn 5, the explicit system exploit (*"Execute a privilege
escalation\..."*) registers a lower absolute residual norm ($2.0650$)
than Turn 3 ($2.8794$). This behavior is caused by lexical overlap
within the base embedder. Because words like *"Execute"* and *"exploit"*
frequently co-occur within the computer science text files used to train
base transformer layers, the raw embedder naturally draws these tokens
closer to the safe programming subspace. However, because HGFW tracks
the sequential state velocity turn-over-turn rather than treating
prompts in isolation, the persistent upward vector pressure still yields
a growth delta of **+0.6037**---safely exceeding the $0.50$ margin and
maintaining containment.

# Threat Model & Adversarial Capabilities Analysis

To evaluate the cryptographic and structural security of the HGFW
architecture, we formalize a multi-tiered threat model analyzing
systemic resilience against varied adversary postures.

## Adversarial Objectives and Assumptions

The adversary's objective is defined as executing an arbitrary code
injection, jailbreak, or privilege escalation payload through the
inference interface. We assume the baseline model layers are fully
frozen, and security must be enforced completely via the HGFW
preprocessing layer. State preservation ($tau_{session}$ and
$prev_{residual}$) is assumed to be externalized to an immutable
distributed database cache (e.g., Redis) mapped to a combined
cryptographic hardware and API token fingerprint, preventing trivial
session-wipe evasion.

## Black-Box Query Interleaving (Token Flood Evasion)

In an interleaved black-box attack vector, the adversary attempts to
pollute or flood the firewall's memory states. They inject an anomalous
payload $x_{bad}$, followed immediately by an engineered sequence of $N$
high-confidence benign payloads $x_{good}$ to force threshold cooldown
or dilute moving-average statistics.\
\
**HGFW Resilience:** Because HGFW avoids sliding-window averages and
utilizes a signed step-difference loop ($|e_t| - |e_{t-1}|$), this
vector fails. When transitioning from a clean query back to a malicious
probe, the vector must mathematically step from the whitelisted space
($|e| \approx 0.0$) back to the anomalous region ($|e| > 2.0$). This
sharp directional turn-over-turn jump violently triggers Gate A
(Structural Variance Tracking) on the first step of re-exploitation,
completely blinding distributed automated botnet fuzzing pipelines.

## White-Box Gradient-Based Steering (Orthogonal Trajectory Matching)

Under a worst-case white-box posture, the adversary is assumed to
possess full access to the target LLM's architecture, its tokenizer
parameters, and the explicit basis weights of the HGFW projection matrix
$P_V$. The attacker can leverage gradient descent (e.g., AutoDAN or
Greedy Coordinate Gradient methods) to optimize a adversarial suffix or
token combination. The mathematical objective of this gradient attack is
to construct a text payload $x^$ that simultaneously satisfies two
constraints:

1.  It contains the required logical exploit tokens to hijack the frozen
    downstream LLM parameters.

2.  Its high-dimensional vector representation satisfies $P_V x^ = x^*$,
    yielding an orthogonal residual component of $|e| = 0$, thereby
    rendering it invisible to the firewall.

**HGFW Resilience:** This mathematical formulation creates a structural
paradox for the attacker. The projection matrix $P_V$ is constructed via
SVD to encapsulate *only* the coordinates of authorized domain behavior.
For an optimized token sequence to achieve $|e| = 0$, its vector
representation must sit entirely within the whitelisted subspace
manifold ($V_k$).If a vector lives completely within $V_k$, it
structurally lacks the mathematical dimensions, orthogonal directional
energy, and token configurations required to steer the downstream LLM's
attention layers out of its authorized operational domain. The adversary
can perfectly match the whitelist coordinates, but doing so
mathematically strips the prompt of the geometric adjustments required
to execute the exploit---rendering the attack completely inert.

# Conclusion & Future Work

HGFW moves deep learning defense away from text-based guessing and
anchors it to rigid linear algebra. Future implementations will scale
the baseline SVD training layer across multi-thousand token public
safety corpuses (such as ShareGPT) to establish real-world
false-positive baselines and mathematically secure the operational
boundaries of enterprise generative AI pipelines.
